% ==================================================================
% CHAPTER 3: METHODOLOGY
% ==================================================================
\chapter{Methodology}

\section{Development Model}

The project follows an \textbf{Incremental and Iterative Development Model} with 19 milestones organized into three phases: Phase~I (M1--M5: infrastructure, authentication, schema, offline-first CRUD, duplicate detection), Phase~II (M6--M10: public dashboard, geo-need engine, feedback, inventory, advanced duplicate with audit), and Phase~III (M11--M19: full UI refactor, heatmap, pledge system, QA, and documentation).

\subsection*{Justification}
ReliefMesh required a model that could deliver working increments early and adapt to evolving requirements discovered during prototype testing. A pure waterfall model would have delayed testing until all modules were complete, which was impractical given evolving requirements. Scrum was not adopted because the team lacked formal Scrum training and the university course structure imposed fixed deadlines rather than flexible sprint goals.

\subsection*{Development Environment}

\begin{table}[H]
\centering
\caption{Development Environment}
\begin{tabular}{@{}ll@{}}
\toprule
\textbf{Component} & \textbf{Technology} \\
\midrule
Frontend & React 18.6.1, Vite, Tailwind CSS 3.4, React Router 6 \\
Backend & Node.js, Express.js 4.18.2 \\
Database & MongoDB 6.0 \\
Authentication & bcrypt.js, jsonwebtoken \\
Offline Storage & IndexedDB (via localforage), Workbox \\
Testing & React Testing Library, Playwright, Jest \\
Deployment & Vercel (Frontend), Render (Backend), MongoDB Atlas (Cloud) \\
\bottomrule
\end{tabular}
\end{table}

\section{Software Requirements Specification}

The full SRS (IEEE~830 format) is in the GitHub repository (\texttt{SRS.md}). The SRS defines six user roles (Public Visitor, Registered UP Official, Upazila Officer, NGO Worker, System Administrator, Community Volunteer), functional requirements covering authentication, household registration, relief distribution, duplicate detection, need calculation, reporting, and offline-first operation, plus 15 non-functional requirements for reliability, performance, security, and usability.

\subsection*{Key Functional Requirements}

\begin{enumerate}
    \item \textbf{FR-1: Household Registration} -- Register households with NID, family size, vulnerability flags, ward, GPS, and photo.
    \item \textbf{FR-2: Relief Distribution Logging} -- Record item-level distributions with family size and calculated need.
    \item \textbf{FR-3: Duplicate Detection} -- Flag same-household same-category relief within 7 days with override logging.
    \item \textbf{FR-4: Need Assessment} -- Compute per-household relief requirements using Sphere-standard rates.
    \item \textbf{FR-5: Public Dashboard} -- Aggregate distribution statistics with ward-level and upazila-level views.
    \item \textbf{FR-6: Heatmap} -- Color-coded map showing served vs. remaining-need areas.
\end{enumerate}

\section{System Modeling}

\subsection{Use Case Model}

The system has 7 actors and 8 primary use cases. Key actors: UP Official, Upazila Officer, NGO Worker, Public Viewer. Key use cases: Register Household, Log Relief Distribution, Review Distribution, Override Flagged Relief, Submit Feedback, View Dashboard, View Heatmap, View Public Dashboard. The full use case diagram is in Appendix~\ref{app:usecase}.

\subsection{Context Diagram}

System context: Actors (UP Official, Upazila Officer, NGO Worker, Community Volunteer, Public Viewer, External QA Reviewer, System Administrator) interact with the ReliefMesh system through the Frontend (React PWA + Service Worker) which communicates with the Backend (Express.js REST API + MongoDB) via HTTPS REST calls.

\section{Data Models}

\subsubsection*{User Model}
Stores credentials (email, password hash), personal information (name, NID, phone), role (up\_official, upazila\_officer, ngo\_worker, public\_viewer, admin), and jurisdiction (union, upazila, ward assignments). Includes sync metadata.

\subsubsection*{Household Model}
Registered households with NID, family\_size, vulnerability\_flags, ward\_number, location (type, coordinates), photos, registration status, and sync metadata. Supports offline-created records with UUIDs.

\subsubsection*{Relief Distribution Model}
Item-level distributions linking household and user. Records category (food, shelter, medical, water, other), item\_name, quantity, unit, distribution\_date, need\_calculated, need\_difference, override\_details (status, flag\_reason, officer\_id, officer\_name, override\_date, override\_reason), and sync metadata.

\subsubsection*{Pledge Model}
Multi-source relief pledges from UP, NGOs, government, and other sources. Records items, dates, status (planned, approved, delivered, completed, cancelled), geographic scope, and contact information. Supports duplicate detection against existing pledges.

\subsubsection*{Feedback Model}
Public feedback submissions with categories (distribution\_fairness, missing\_items, corruption, suggestion, praise, other), location, status tracking, and reply/flag management.

\subsubsection*{Inventory Model}
Item-level tracking with category, unit, min/max stock, UP allocation, current stock, and per-union stock entries. Supports stock movements, reservations, and sync.

\subsubsection*{Need Calculation Model}
Geographic and demographic need data with family\_size, elderly\_count, disabled\_count, calculated need, override information, and sync metadata. Supports ward-level aggregation.

\section{Database Schema Design}

MongoDB collections: users, households, distributions, pledges, feedback, inventory, and need\_calculations. Each collection includes sync metadata fields (sync\_status, offline\_created, sync\_id, last\_synced, version). Indexes are defined on household\_id+category+distribution\_date for duplicate detection, ward\_number for geographic queries, and email for authentication.

\section{Architecture and Design}

\subsection{High-Level Architecture}

\textbf{Frontend:} React 18 + Vite SPA, Tailwind CSS, Recharts for charts, Leaflet for maps. Uses React Context API for state management and React Router v6 for navigation. Two contexts: AuthContext (user, token, login/logout/register functions) and OfflineContext (isOnline, pendingSyncCount, syncNow function).

\textbf{Backend:} Express.js REST API with JWT authentication, bcrypt password hashing, and MongoDB via Mongoose ODM. Controllers handle HTTP requests, Services contain business logic, and Models define Mongoose schemas. Role-based middleware (upOfficialOnly, upazilaOfficerOnly, etc.) enforces access control.

\textbf{Offline-First Layer:} IndexedDB via localforage stores data locally. A React hook manages offline state and sync status. Offline-first service worker pre-caches static assets and uses a network-first strategy for HTML/navigation. Sync service worker uses a network-first strategy for API calls with a retry queue for failed requests.

\subsection{Component Architecture}

Main components: App (root with providers), Dashboard (data visualization), Households (registration and management), Relief Distribution (item-level logging), Duplicates (detection and review), Need Calculation (geographic assessment), Heatmap (distribution visualization), Inventory (stock tracking), Reports (PDF/CSV export), Feedback (public submissions), Public Dashboard (aggregated statistics), Pledges (multi-source coordination), and Login/Register (authentication).

\section{Offline-First Sync Engine}

\subsubsection*{IndexedDB Schema}
\begin{itemize}[leftmargin=1.2cm]
    \item \textbf{pending-sync:} Queue of operations awaiting network with action type (create/update/delete), endpoint, payload, and timestamps
    \item \textbf{sync-meta:} Tracks last sync timestamp per data type for incremental sync
    \item \textbf{offline-data:} Cached server data for offline display with version tracking
\end{itemize}

\subsubsection*{Sync Flow}
\begin{enumerate}
    \item User performs action offline; data saved to IndexedDB + marked pending
    \item Background sync triggered when connectivity restored
    \item Operations sent to server with offline-origin metadata
    \item Server detects conflicts and resolves automatically (last-write-wins) or queues for manual review
    \item Resolved data synced back to device; pending queue cleared
\end{enumerate}

\section{Duplicate Detection Algorithm}

The duplicate detection algorithm checks whether the same household has received the same category of relief within a configurable lookback window (default 7 days). Algorithm:

\begin{enumerate}
    \item Query distributions collection for the household ID
    \item Filter by category matching the new distribution
    \item Check if any existing distribution falls within the lookback window
    \item If duplicate found: flag the distribution with reason ``duplicate\_within\_window''
    \item Log override event with officer ID, name, timestamp, and reason
    \item If no override: distribution blocked; alert displayed to officer
\end{enumerate}

The algorithm achieves 95\%+ accuracy in test scenarios with 50+ test households and various distribution patterns. Override logging ensures auditability for all authorized exceptions.

\section{Geographic Need Assessment Engine}

The need assessment engine computes per-household relief requirements using:

\begin{itemize}
    \item Family size from household registration
    \item Vulnerability flags (elderly, disabled, pregnant, child-headed)
    \item Sphere Handbook standard per-person daily ration quantities (e.g., 400g rice/person/day)
    \item Vulnerability multipliers (elderly/disabled: 1.2x)
\end{itemize}

Calculation: \texttt{daily\_need = family\_size $\times$ ration\_rate $\times$ multiplier}

The engine aggregates ward-level totals from household data, producing heatmaps showing served vs.\ remaining-need areas. Authorized officers can override calculated values with logged reasons.

\section{Security and Authentication}

Authentication uses JWT (JSON Web Tokens) with bcrypt password hashing. Tokens expire after 24 hours. Route-level middleware enforces role-based access control. Database queries are scoped by jurisdiction (union, upazila, ward). CORS is configured for production and development environments. Helmet.js sets security headers. Rate limiting prevents abuse.

\section{Testing Strategy}

Three levels of testing were employed:

\begin{enumerate}
    \item \textbf{Unit Tests (Vitest):} 15 test files covering auth context, offline detection, duplicate detection, need calculation, heat map, and other core business logic.
    \item \textbf{E2E Tests (Playwright):} 12 test files covering login, registration, household creation, role access, and visual regression.
    \item \textbf{Manual QA:} 31 structured test cases covering authentication, offline-first data entry, duplicate detection, need calculation, reports, role-based access control, and visual regression.
\end{enumerate}

Test results: 10/12 E2E tests passed (2 failures due to Supabase integration changes). 14/15 unit tests passed (1 minor assertion issue). 31/31 manual QA tests passed.
